\section{未来展望与结论}

\begin{frame}{未来切入点：按“先可信、再自主、后闭环”推进}
\scriptsize
\begin{tabularx}{\textwidth}{p{0.07\textwidth} p{0.24\textwidth} X p{0.13\textwidth} p{0.13\textwidth}}
\toprule
序 & 切入点 & 直接收益 & 价值 & 难度 \\
\midrule
1 & 测量系统与数据契约 & precision/coverage 可信，消除 summary 误读 & 极高 & 中 \\
2 & 可恢复 full campaign & 初始语料按有效证据合同交付，长任务可 resume & 极高 & 中高 \\
3 & source-aware policy & 固定预算下减少 source skew 与低价值调用 & 高 & 高 \\
4 & active relevance & 用少量人工复核压缩 uncertain，改善长尾 & 高 & 中 \\
5 & adversarial/resilience & 证明恶意源、429/5xx、SDK/Mongo 故障下仍安全 & 高 & 中 \\
6 & global KG closed loop & 让结构空洞反向产生 query，形成动态 LLM 安全图 & 极高 & 很高 \\
\bottomrule
\end{tabularx}

\vspace{0.7em}
\begin{block}{推荐资源配置}
前两项是所有后续研究的共同地基，应占首个里程碑的大部分投入；bandit、LLM 自主 query 与 KG 闭环只有在 measurement contract 稳定后才有可信收益。
\end{block}
\end{frame}

\begin{frame}{分阶段路线图：工程与论文共用同一证据资产}
\centering
\begin{tikzpicture}[
  n/.style={draw,rounded corners,minimum width=3.0cm,minimum height=1.4cm,align=center,font=\scriptsize,thick},
  a/.style={-{Latex[length=1.7mm]},thick,Navy},node distance=0.45cm
]
\node[n,fill=SoftBlue] (p1) {0-30 天\\MSA/rubric\\统一 gap-stop schema\\frozen fixtures};
\node[n,fill=SoftGreen,right=of p1] (p2) {31-60 天\\full campaign MVP\\checkpoint/resume\\Control dashboard};
\node[n,fill=SoftAmber,right=of p2] (p3) {61-90 天\\同预算 replay/DOE\\active relevance\\adversarial suite};
\node[n,fill=SoftRed,right=of p3] (p4) {3-12 个月\\contextual policy\\benchmark paper\\global KG loop};
\draw[a] (p1)--(p2);\draw[a] (p2)--(p3);\draw[a] (p3)--(p4);
\end{tikzpicture}

\vspace{1em}
\begin{block}{每阶段只有一个共同 gate}
产物必须同时服务工程发布和学术复现：同一个 task manifest、source snapshot、annotation、trace、config 与统计脚本，不维护两套“展示数据”。
\end{block}
\end{frame}

\begin{frame}{结论：当前系统的上限由算法决定，下限由测量与控制决定}
\begin{columns}[T,onlytextwidth]
\column{0.48\textwidth}
\begin{exampleblock}{源码证明的优势}
混合 controller、白名单工具、确定性复核、有效证据、UCB、checkpoint、playbook、Mongo provenance 与 trace 已形成一条具有研究潜力的完整骨架。
\end{exampleblock}
\column{0.48\textwidth}
\begin{alertblock}{运行证明的现实}
full/legacy 路径、来源偏斜、uncertain 过高、终止字段歧义使“完成采集”仍缺少可信合同；latest run 是这个矛盾的直接证据。
\end{alertblock}
\end{columns}

\vspace{1em}
\begin{center}
\begin{minipage}{0.86\textwidth}
\Large\centering
先用 DMAIC 把\textcolor{Blue}{测量、交付与控制}做实，
再把\textcolor{Purple}{主动搜索、风险停止与 KG 反馈}做成可证伪的学术贡献。
\end{minipage}
\end{center}
\end{frame}

\begin{frame}{交付物与复现入口}
\scriptsize
\begin{tabularx}{\textwidth}{p{0.30\textwidth} X}
\toprule
文件 & 用途 \\
\midrule
主报告 \,\pathcode{.tex/.pdf} & Beamer 主报告与编译产物 \\
\pathcode{slides/*.tex} & 分章节可维护源码 \\
\pathcode{code_evidence_matrix.md} & 报告结论到源码位置的完整索引 \\
\pathcode{runtime_snapshot.json} & Mongo、latest run、trace、tests 的统一快照 \\
\pathcode{literature_review.md} & 检索协议、文献矩阵、综合结论与注释书目 \\
\pathcode{references.bib} & DOI/URL 核验的 BibLaTeX 文献库（实际文件名见 README） \\
\pathcode{build_report.ps1} & XeLaTeX + Biber 可复现构建 \\
\bottomrule
\end{tabularx}

\vspace{0.7em}
\begin{block}{复核建议}
先看运行结果与“容易误读”页，再按 code evidence matrix 回到代码；需要做改进时，从 DMAIC Measure tollgate 开始，不从替换模型开始。
\end{block}
\end{frame}

\section{参考文献}

\begin{frame}[allowframebreaks]{参考文献}
\tiny
\printbibliography[heading=none]
\end{frame}
